\documentclass[openany, a4paper, 12pt]{book}

% ==========================================
% main2.tex — 动机社会重构版（PREVIEW SCAFFOLD）
% ==========================================
% ⚠️ 本文件是重构版入口，与 main.tex 完全独立。
%    main.tex 与 part1_core_logic/ 等原章节文件一律不改动。
%    本文件是 subfiles 宿主，chapter 文件用 \documentclass[../main2.tex]{subfiles}
% ==========================================
% 风格依据：补充材料/part3_AI执行指导.md + skills/motivation-society-style/SKILL.md
% 架构依据：notes/00-architecture.md
% ==========================================

% ==========================================
% 1. 基础宏包配置（与 main.tex 保持一致，macOS 已验证）
% ==========================================
\usepackage[UTF8,fontset=none]{ctex}     % 中文支持（手动指定字体，避免 TinyTeX mac fontset 缺失）
\setCJKmainfont{Songti SC}
\setCJKsansfont{Songti SC}
\setCJKmonofont{Menlo}
\setmainfont{Times New Roman}
\setsansfont{Helvetica}
\setmonofont{Menlo}
\usepackage{unicode-math}
\setmathfont{STIX Two Math}
\usepackage{geometry}       % 页面设置
\geometry{left=1in, right=1in, top=1in, bottom=1in}
\usepackage{fancyhdr}       % 页眉页脚
\usepackage{hyperref}       % 超链接目录
\hypersetup{colorlinks=true, linkcolor=blue, filecolor=magenta, urlcolor=cyan}
\usepackage{subfiles}       % 子文件支持

% ==========================================
% 2. 绘图与框图工具
% ==========================================
\usepackage{tikz}
\usetikzlibrary{shapes.geometric, arrows.meta, positioning, shadows, calc, fit, trees}

\usepackage[most]{tcolorbox}
\tcbuselibrary{skins, breakable, theorems}

% 阅读笔记框
\newtcolorbox{readingnote}{
	colback=gray!5,
	colframe=gray!40,
	boxrule=0.4pt,
	arc=2mm,
	left=6pt,
	right=6pt,
	top=6pt,
	bottom=6pt
}

\newcommand{\xr}{\xrightarrow}

% ==========================================
% 数学附录 ↔ 正文 链接接口
% ==========================================
\newcommand{\mathlink}[2]{%
	\par\noindent
	\textcolor{blue}{\small$\rightarrow$ \hyperref[#1]{#2}}
	\par
}

% ==========================================
% 低频高信息密度强调
% ==========================================
\newcommand{\lowfreq}[1]{\textbf{\textcolor{red!70!black}{#1}}}

% ==========================================
% 3. 自定义环境定义
% ==========================================

% 逻辑地图
\newenvironment{logicmap}{
	\begin{tcolorbox}[
		colback=blue!5!white,
		colframe=blue!50!black,
		title=\textbf{逻辑地图 (Logic Map)},
		fonttitle=\bfseries,
		sharp corners=south,
		boxrule=0.5mm,
		]
	\centering
}{
	\end{tcolorbox}
}

% 本章动机
\newenvironment{motivation}{
	\vspace{0.5em}
	\begin{tcolorbox}[
		colback=gray!10,
		colframe=gray!60,
		title=\textit{本章动机 (Motivation)},
		fonttitle=\bfseries,
		arc=0mm,
		leftrule=3mm
		]
		\small \itshape
}{
	\end{tcolorbox}
	\vspace{1em}
}

% 数学推导盒子
\newtcolorbox{mathbox}[1]{
	enhanced,
	colback=gray!5,
	colframe=gray!60!black,
	title=\textbf{数学建模：#1},
	fonttitle=\small\bfseries\sffamily,
	leftrule=2mm,
	breakable
}

\newtcolorbox{definitionbox}{
	colback=blue!3!white,
	colframe=blue!60!black,
	fonttitle=\bfseries,
	breakable,
	enhanced,
	boxrule=0.8pt,
	arc=2mm
}
\newtcolorbox{questionbox}{
	colback=red!3!white,
	colframe=red!60!black,
	breakable,
	enhanced,
	boxrule=0.8pt,
	arc=2mm
}
\newtcolorbox{examplebox}{
	colback=green!3!white,
	colframe=green!50!black,
	breakable,
	enhanced,
	boxrule=0.8pt,
	arc=2mm
}
\newtcolorbox{warningbox}{
	colback=orange!3!white,
	colframe=orange!60!black,
	breakable,
	enhanced,
	boxrule=0.8pt,
	arc=2mm
}

\newenvironment{reader_anticipation}{
	\begin{tcolorbox}[colback=yellow!3!white,colframe=yellow!50!black,title=\textit{读者可能还会问},breakable]
	\begin{itemize}
}{
	\end{itemize}
	\end{tcolorbox}
}

\usepackage{booktabs}
\usepackage{amsmath,amssymb}

% ==========================================
% 4. 书籍信息
% ==========================================
\title{\textbf{谁在推动未来}\\ \large ——从「发生了什么」到「谁在起作用」的重构 (重构版 v2)}
\author{D·declaim}
\date{\today}

% ==========================================
% 5. 正文开始
% ==========================================
\begin{document}

\maketitle

\frontmatter
\subfile{frontmatter/preface_v2}
\tableofcontents

\mainmatter

% =================================================================
% 导论：为什么算得最准的人，预测得最离谱
% =================================================================
\chapter*{导论：为什么算得最准的人，预测得最离谱}
\addcontentsline{toc}{chapter}{导论：为什么算得最准的人，预测得最离谱}
\subfile{part1_motivation_goal/introduction_v2}

% =================================================================
% PART I: 驱动未来的引擎
% =================================================================
\part{驱动未来的引擎：动机与目标 (The Engine)}
\subfile{part1_motivation_goal/chapter1_motivation}
\subfile{part1_motivation_goal/chapter2_goal}

% =================================================================
% PART II: 谁是真正的推手
%   v2 变化：原「演化状态空间」整部删除。因果推断从「预测未来」内部长出。
% =================================================================
\part{谁是真正的推手：预测与因果 (Prediction and Causation)}
\subfile{part2_prediction_causality/chapter6_prediction_cause}
\subfile{part2_prediction_causality/chapter7_correlation_causation}

% =================================================================
% PART III: 法庭上的时光机
%   v2 变化：法律从「Part IV 的一个通道」升格为「方法最早的成熟应用」
% =================================================================
\part{法庭上的时光机：多因分账 (The Courtroom)}
\subfile{part3_courtroom/chapter8_butfor_test}
\subfile{part3_courtroom/chapter9_multiple_causes}

% =================================================================
% PART IV: 社会的四大齿轮
% =================================================================
\part{社会的四大齿轮：观念、经济、权力、法律 (Four Gears)}
\subfile{part4_gears/chapter10_channel_generator}
\subfile{part4_gears/chapter11_meme}
\subfile{part4_gears/chapter12_economy}
\subfile{part4_gears/chapter13_politics}
\subfile{part4_gears/chapter14_law}

% =================================================================
% PART V: 回到循环
% =================================================================
\part{回到循环：预测如何改变预测 (The Loop)}
\subfile{part5_loop/chapter15_forecast}
\subfile{part5_loop/chapter16_closure}

% =================================================================
% 数学附录
% =================================================================
\appendix
\subfile{appendix/appendix_v2_math}

\backmatter
\subfile{backmatter/acknowledgements_v2}

\end{document}
